% Misconfigured by Default — Springer Nature (sn-jnl) skeleton for EMSE submission.
% Reference style: sn-basic (author-year), per EMSE's name-year citation rule.
\documentclass[pdflatex,sn-basic]{sn-jnl}

\usepackage{amsmath,amssymb,amsfonts}
\usepackage{textcomp}
\usepackage{booktabs}
\usepackage{longtable}
\usepackage{array}
\usepackage{calc}
\usepackage{etoolbox}
\providecommand{\tightlist}{\setlength{\itemsep}{0pt}\setlength{\parskip}{0pt}}
\AtBeginEnvironment{longtable}{\footnotesize\setlength{\tabcolsep}{4pt}}

\title[Misconfigured by Default]{Misconfigured by Default: The Prevalence, Validity, and Tool Coverage of Kubernetes Misconfigurations in the Wild}

\author*[1]{\fnm{Girish} \sur{Mallya}}\email{girishm60@gmail.com}

\affil*[1]{\orgname{Independent Researcher}, \orgaddress{\city{Bellevue}, \state{WA}, \country{USA}}}

\abstract{Kubernetes manifests are the de facto interface to production
infrastructure, with permissive defaults: a workload specifying nothing
runs with an automounted token, no resource limits, and no sandboxing
profile. We measure how 151 Helm charts and 60 public GitHub
repositories (16,243 documents, 2,811 workloads) configure these knobs,
using a 25-rule engine.

We validated the engine on a 457-finding stratified sample: AI-generated
pre-labels were author-checked --- the 92 protocol-sensitive rows
individually adjudicated, a blinded 30-row subset labeled blind
(author--AI agreement: $\kappa$ = 0.526 {[}0.17, 0.84{]}) --- and the rest
retain pre-labels accepted in the verification pass. Population-weighted
precision is 0.699 {[}0.62, 0.84{]} (cluster bootstrap by source):
$\approx$19,200 {[}17,100, 23,100{]} of 27,460 raw findings are estimated
genuine. Precision is rule-dependent: unambiguous absence checks
(seccomp, privilege escalation, root prevention) are essentially exact,
while context-sensitive rules (committed secrets, host namespaces,
probes) are mostly benign-by-design cases.

The defaults dominate: tokens automount in 88.6\% of production GitHub
workloads, CPU limits --- a contested hardening practice --- are missing
in 75.3\%; 1.4\% run a privileged container. Helm charts beat raw
manifests on most practices, but by only 1.3--1.7$\times$ in workload-level
prevalence.

We also audit the auditors: on identical input, four widely used
scanners cover at most 16 of the 25 practices, diverging up to 1.4$\times$ on
the same practice largely through counting semantics. Scored against our
verified labels (Helm arm, n = 126), they effectively detect 52--75\% of
the engine's verified-genuine findings and flag 86--95\% of
verified-benign findings where a check exists, under our contextual
protocol.

static analysis tools; measurement validity}

\keywords{Kubernetes, configuration errors, empirical study, misconfiguration,}

\begin{document}

\maketitle

\hypertarget{introduction}{%
\subsection{1. Introduction}\label{introduction}}

Almost every production system now runs on Kubernetes, and almost every
Kubernetes workload is described by a YAML manifest that its author
wrote once and rarely revisits. The platform's configuration surface is
large, and its defaults were chosen for convenience rather than safety.
A container that says nothing about security runs as whatever user its
image specifies, with a service-account token mounted into its
filesystem, no CPU or memory limits, no seccomp profile, and privilege
escalation allowed. Each of these is a documented hardening gap; in
aggregate, they define the actual security posture of the cloud-native
ecosystem.

Two questions follow. First, \emph{descriptive}: how are real manifests
actually configured --- which hardening practices are routine, which are
rare, and how much of the apparent risk is an artifact of counting test
fixtures and examples as if they were production code? Second,
\emph{instrumental}: the ecosystem's answer to this problem is a family
of static auditing tools, but those tools are themselves rarely audited.
If four widely used scanners are run on identical input, do they cover
the same practices, and do their counts agree?

Both questions are measurement questions, and both are easy to answer
badly. Manifest corpora scraped from public sources mix production code
with tutorials, test fixtures, and rendered templates; naive counting
silently grades homework as if it were infrastructure. Scanner
comparisons, meanwhile, tend to report raw finding counts without
establishing whether any tool's findings are \emph{correct} --- and our
own validation shows that for several common rules, the majority of
flagged findings are benign in context.

This paper makes four contributions:

\begin{enumerate}
\def\labelenumi{\arabic{enumi}.}
\tightlist
\item
  \textbf{A two-arm prevalence measurement} of 25 misconfiguration
  practices across 16,243 documents (2,811 workloads), with source
  tagging that separates production manifests from examples and
  fixtures, cluster-level prevalence that corrects for duplicated
  content, and a Helm-vs-GitHub comparison that corrects a best-chart
  bias in the obvious sampling strategy (Section 4.1, 4.2).
\item
  \textbf{A validated rule engine.} Every prevalence number is paired
  with a per-rule precision estimate from a 457-finding stratified
  sample, judged under a written labeling protocol whose evolution
  (three versions), blinded author--AI agreement check, and adjudication
  amendments are documented and released. Precision, not raw counts, is
  the paper's unit of evidence (Section 4.3).
\item
  \textbf{A revised taxonomy} of Kubernetes misconfiguration that
  distinguishes misconfigurations proper from hardening gaps and
  manifest rot, and classifies false positives into reusable context
  classes (infrastructure agents, credential-free placeholders, standard
  operator patterns) (Section 4.3).
\item
  \textbf{An audit of four production scanners on identical input},
  showing that coverage tops out at \textasciitilde16 of 25 practices,
  that agreement is high only where misconfiguration is explicit, and
  that headline count divergences (up to \textasciitilde1.4$\times$ on the same
  practice) are explained by counting semantics --- pod- vs
  container-level units and merged checks; and that, scored against
  verified ground truth, the tools' real deficits are coverage and
  context-sensitivity, not detection on the practices they do check
  (Section 4.4).
\end{enumerate}

The headline result is that in this corpus the risk is concentrated in
\emph{defaults}, not in explicit misbehavior. Privileged containers
appear in 1.4\% of production workloads; automounted service-account
tokens appear in 88.6\%. The pattern suggests security effort --- human
and tooling alike --- is calibrated to the rare, explicit cases and
largely blind to the common, default ones.

\hypertarget{background-and-research-questions}{%
\subsection{2. Background and research
questions}\label{background-and-research-questions}}

A Kubernetes workload manifest declares, among other things, a pod's
security context (user, capabilities, seccomp and AppArmor profiles,
privilege escalation), its containers' resource requests and limits,
image references (tag and digest pinning), health probes,
service-account binding, and --- in adjacent objects --- RBAC roles and
Secrets. Nearly every field is optional, and the platform supplies a
permissive default when a field is absent. Misconfiguration, for this
study, is a manifest state that a documented hardening practice (the CIS
Kubernetes Benchmark, the NSA/CISA hardening guide, and the scanners'
own rule catalogues) identifies as a gap. Section 3.2 defines the 25
rules operationally; Section 3.4 describes how we establish that the
rules measure what they name.

We ask four research questions:

\begin{itemize}
\tightlist
\item
  \textbf{RQ1 (Prevalence).} How common is each class of
  misconfiguration in real-world manifests, once production code is
  separated from examples and fixtures and duplicated content is
  accounted for?
\item
  \textbf{RQ2 (Validity and taxonomy).} With what precision does a
  rule-based engine detect each practice --- i.e., which of the 25
  measurements can be taken at face value and which require contextual
  judgment --- and what taxonomy organizes the true and false positives?
\item
  \textbf{RQ3 (Tool coverage).} On identical input, how much of the
  observed taxonomy do the four widely used scanners cover, and do their
  counts agree? Where they diverge, is the divergence detection or
  bookkeeping?
\item
  \textbf{RQ4 (Curation effect).} Do rendered Helm charts --- the
  ecosystem's curated, templated distribution channel --- differ from
  manifests committed directly to Git repositories, and by how much once
  sampling bias is handled?
\end{itemize}

\hypertarget{method}{%
\subsection{3. Method}\label{method}}

\hypertarget{corpus-construction}{%
\subsubsection{3.1 Corpus construction}\label{corpus-construction}}

The corpus has two arms, collected in October 2026 from public sources
only.

\textbf{Helm arm.} We rendered the 151 most popular charts on Artifact
Hub with default values (\texttt{helm\ template}), yielding 2,616
documents and 402 workloads. Default values are what an operator gets
when they install a chart without tuning it --- the correct object of
study for a question about defaults --- but a top-charts sample
over-represents well-maintained projects; Section 4.2 quantifies and
corrects for this \emph{best-chart bias} rather than ignoring it.

\textbf{GitHub arm.} We mined 60 public repositories (seed projects plus
search-discovered manifest collections), yielding 13,627 documents and
2,409 workloads. Snapshots are pinned by commit SHA
(\texttt{results/github\_commit\_shas.json}). Every file was
source-tagged as \emph{production}, \emph{example}, or \emph{fixture} by
path and repository role, so prevalence can be reported for production
manifests alone; an addendum to the labeling protocol (Section 3.4)
resolves tag/content conflicts in favor of content (e.g., a
\texttt{testdata/} tree inside a production-tagged repository counts as
fixture).

The corpus itself is not redistributed (third-party manifests carry
mixed licenses); the collection scripts and per-source logs are the
corpus manifest, and running the collectors rebuilds it.

\hypertarget{the-rule-engine}{%
\subsubsection{3.2 The rule engine}\label{the-rule-engine}}

We implemented a 25-rule engine over parsed YAML (Deployments,
StatefulSets, DaemonSets, Jobs, CronJobs, Pods, and the RBAC/Secret
objects adjacent to them). Rules span six families: \textbf{images}
(untagged, \texttt{:latest}, no digest pinning), \textbf{resources}
(missing CPU/memory limits, missing requests), \textbf{reliability}
(missing liveness/readiness probes), \textbf{security context}
(privileged, root not prevented, explicit root, privilege escalation not
disabled, dangerous or any added capabilities, missing seccomp),
\textbf{host access} (host network/PID/IPC, automounted SA token),
\textbf{RBAC and secrets} (wildcard verbs/resources, cluster-admin
binding, committed Secret objects, literal secrets in environment
variables), and \textbf{API currency} (deprecated APIs --- reported as
manifest rot, not misconfiguration; see Section 4.3). Engine v3
incorporates three correctness fixes over v2 that matter for the
results: the SA-token rule joins the ServiceAccount object's own
automount flag; init containers are excluded from probe rules (a third
of v2's Helm probe findings were on init containers); and
\texttt{seccompProfile.type=Unconfined} counts as no profile. Engine v3
reports 27,460 findings over the full corpus. The full rule definitions
ship with the replication package.

\hypertarget{tool-audit-setup-rq3}{%
\subsubsection{3.3 Tool audit setup (RQ3)}\label{tool-audit-setup-rq3}}

We ran kube-linter, kubeaudit, Polaris, and Checkov --- the four most
widely adopted open-source manifest auditors --- on the \emph{identical}
Helm-arm input (402 workloads, rendered with default values), each with
its default configuration. Default configuration is the treatment of
interest: it is what a new adopter runs, and, as Section 4.4 shows,
several relevant checks are not in it. We mapped each tool's checks to
our 25 rules by hand from the tools' output on this input, recording
partial coverage (a check that covers a rule with narrower semantics) as
one half in coverage totals. Counts were reconciled rule by rule; where
a tool merges two of our rules into one check (e.g., untagged and
\texttt{:latest}), we report the merged figure rather than forcing an
attribution. Separately, we score each tool against ground truth: the
126 Helm-arm rows of the validation sample (Section 3.4) carry
author-verified labels, so for every sampled finding we can ask whether
the tool reported the corresponding practice on that object --- yielding
detection rates on verified true positives and flagging rates on
verified false positives, per tool (full method and per-row data in
\texttt{results/tool\_accuracy.md}).

\hypertarget{validation-labeling-protocol-reliability-adjudication}{%
\subsubsection{3.4 Validation: labeling protocol, reliability,
adjudication}\label{validation-labeling-protocol-reliability-adjudication}}

Prevalence without precision is numerology, so we validated the engine
on a 457-finding stratified sample: all findings of the two rarest
high-severity rules, stratified draws from the rest, and
\textasciitilde10\% hidden duplicate rows. A sample row is one finding
with an \emph{evidence packet} (the flagged object, its security
context, and its source context).

Labels (TP / FP / unsure) were assigned under a written protocol. The
protocol's history is part of the method and is reported, not hidden: an
initial \emph{definitional} standard (v2: a finding is TP iff the
manifest exhibits the property the rule names, in context) was
harmonized with an earlier \emph{exploitability} draft (v1) into a
governing protocol (v3) that keeps the definitional core and adds
explicit contextual amendments --- production scoping for resource/probe
rules, a two-tier treatment of digest pinning, an escalation-path test
for root/privilege-escalation rules, and a credential-material test for
committed secrets. AI pre-labels produced under v2 were treated as an
\emph{initial annotation pass}, never as ground truth.

\textbf{Blinded author--AI agreement.} The author labeled a blinded
30-row subset (spanning all 25 rules) before seeing any pre-labels.
Agreement was 24/30 (80\%), Cohen's $\kappa$ = 0.526 {[}bootstrap 95\% CI 0.17,
0.84{]} against the v2-standard pre-labels. This is agreement between
the author and AI-generated labels, not inter-rater reliability between
two humans; and because the pre-labels follow v2 while the author's
labels follow the emerging v3, $\kappa$ also absorbs the protocol-version
difference. A two-human measurement --- an independent researcher
labeling a 150-row subsample under the frozen v3 protocol, blind to all
existing labels --- is in preparation; the packet ships in the
replication package. We report this number as measured, with its
decomposition, because the decomposition is the finding: four of six
disagreements were resource/probe rows where the governing v3
production-scoping amendment sides with the author against the v2
definitional labels --- a standards difference, not rater noise --- and
the remaining two were genuine judgment calls (a CNI agent's dangerous
capabilities; an operator's wildcard RBAC confined to its own API
groups), which were adjudicated into two further \emph{rule-level}
amendments (an infrastructure-class test requiring a drop-ALL baseline
plus function-required capabilities, and an own-API-group scoping test
for wildcard RBAC), adopted before the verification pass and now part of
the protocol (\texttt{results/labeling\_protocol\_v3\_addendum.md}).
Intra-rater consistency, measured by the hidden duplicates, was 44/44
groups.

\textbf{Verification.} The author then ran a verification pass over the
full sample: the 92 \emph{divergence rows} --- those whose labels can
differ between the v1 and v2 standards --- were adjudicated first, and
the author accepted the v3 recommendations on all 92 without override on
2026-10-27; the remaining rows were reviewed in aggregate (blank =
accept, per protocol). The provenance of the final labels is therefore
mixed, and we state it exactly: of the 451 sample rows that persist in
engine v3, \textbf{111 (25\%) carry labels the author produced or
individually adjudicated} --- the 30 blind-labeled rows and 87
adjudicated divergence rows, with 6 rows in both sets --- while
\textbf{340 (75\%) retain the AI pre-labels} the author accepted during
the pass. The validation dataset is thus AI-annotated and
author-adjudicated, not independently human-labeled; what that does and
does not establish is taken up in Section 6. Rows whose flagged property
cannot be observed in the source (7, e.g.~images injected by kustomize
overlays) are excluded from precision denominators with the count and
reason disclosed.

\hypertarget{analysis}{%
\subsubsection{3.5 Analysis}\label{analysis}}

For each rule we report precision = TP / (TP + FP) over decided rows
with Wilson 95\% confidence intervals, \emph{adjusted prevalence} =
precision $\times$ population findings, and workload-level prevalence (share of
workloads with $\geq$1 finding) per corpus arm and source tag. Because the
validation sample was drawn under engine v2 and the population is engine
v3, we reconciled every sampled row against v3 findings on (arm, source,
kind, name, container, rule): 451 of 457 rows still exist in v3; the
headline figures are computed on that matched frame, and a
distinct-finding sensitivity check (one vote per unique key, collapsing
the intentional duplicates) is reported alongside. Two aggregates of the
per-rule precisions answer different questions and are kept distinct
throughout: the rate across sampled rows (0.633) describes the
\emph{sample}, in which rare rules are deliberately over-represented
because all findings of the rarest rules were taken; the
population-weighted aggregate --- each rule's precision weighted by its
v3 population, equivalently the adjusted finding count over all findings
--- is \textbf{0.699} (19,199 / 27,460) and is the figure that
characterizes the corpus as a whole. Propagating the per-rule binomial
sampling variances, the adjusted intervals are cluster-bootstrapped by
source (resampling repositories and charts, 2,000 replicates;
\texttt{results/cluster\_bootstrap.md}): the adjusted total's 95\%
interval is {[}17,120, 23,071{]} and the weighted precision's is
{[}0.623, 0.840{]}. Row-level intervals that treat findings as
independent are materially narrower and are not used, since findings
within a source are correlated (52.9\% of v3 findings repeat a seen
key). The intervals are also insensitive to how a replicate treats a
rule with no decided rows in the resample --- an event in only 42 of
50,000 rule-replicates: full-sample, pooled-precision, and
drop-and-renormalize fallbacks all yield {[}0.623, 0.840{]}. Every
adjusted prevalence figure in Section 4 is built per rule and is
identical under either aggregate. As a standards sensitivity check, the
same sample scored under the superseded v2 definitional standard (the
pre-labels) gives $\approx$0.73; the gap to 0.633 is the v3 amendments ---
production scoping and the context classes of Section 4.3 ---
reclassifying findings that the purely definitional standard counted as
genuine. Cluster-level prevalence (share of repositories/charts with $\geq$1
finding) addresses the non-independence of findings within a source:
52.9\% of v3 GitHub findings repeat an already-seen (kind, name,
container, rule) key, so document-level rates alone would overstate how
widespread each practice is --- and, as Section 4.1 shows, the practices
are widespread anyway.

\begin{center}\rule{0.5\linewidth}{0.5pt}\end{center}

\hypertarget{results}{%
\subsection{4. Results}\label{results}}

\hypertarget{rq1-prevalence-the-defaults-are-the-finding}{%
\subsubsection{4.1 RQ1 --- Prevalence: the defaults are the
finding}\label{rq1-prevalence-the-defaults-are-the-finding}}

Table 1 reports workload-level prevalence for the headline practices,
per arm, with the GitHub arm restricted to production-tagged manifests.
Three facts stand out.

\textbf{Default-insecure settings are near-universal.} In production
GitHub manifests, service-account tokens are automounted in
\textbf{88.6\%} of workloads and CPU limits are missing in
\textbf{75.3\%}; in rendered Helm charts the figures are 85.1\% and
85.6\%, with \textbf{97.5\%} of Helm workloads running at least one
image that is not digest-pinned (65.3\% for production GitHub).
Restricting to production code barely moves the numbers (corpus-wide
GitHub SA automount is 96.2\%): the risky defaults live in shipped
manifests, not in test files.

\textbf{Explicit dangers are rare.} Privileged containers appear in
\textbf{1.4\%} of production GitHub workloads (5.5\% in Helm charts,
where node-level agents are common), and cluster-level prevalence tells
the same story from the other side: SA automount appears in
\textbf{94\%} of repositories and \textbf{90\%} of charts, missing
seccomp in 94\%/72\%, unpinned digests in 91\%/97\% --- but privileged
containers in only 22\% of repositories and 8\% of charts. The ecosystem
has largely learned not to do the obviously dangerous thing, and has
largely not learned to turn off the defaults.

\textbf{Table 1.} Workload-level prevalence, engine v3 (share of
workloads with $\geq$1 finding).

\begin{longtable}[]{@{}lrrr@{}}
\toprule\noalign{}
Practice & Helm charts & GitHub (production) & GitHub (all) \\
\midrule\noalign{}
\endhead
\bottomrule\noalign{}
\endlastfoot
Image not digest-pinned & 97.5\% & 65.3\% & --- \\
SA token automounted & 85.1\% & 88.6\% & 96.2\% \\
No CPU limit & 85.6\% & 75.3\% & 79.0\% \\
Privileged container & 5.5\% & 1.4\% & --- \\
\end{longtable}

\emph{Full per-rule tables for all 25 rules, both arms and all source
tags, are in the replication package
(\texttt{results/stats\_tagged\_v3.json}).}

\hypertarget{rq4-curation-helps-less-than-it-looks}{%
\subsubsection{4.2 RQ4 --- Curation helps, less than it
looks}\label{rq4-curation-helps-less-than-it-looks}}

Rendered Helm charts are better hardened than raw GitHub manifests, but
the gap is \textbf{1.3--1.7$\times$}, not the 2--3$\times$ suggested by comparing the
most popular charts against unfiltered repository mining. The mechanism
is sampling: a top-151-charts frame over-weights flagship projects with
dedicated maintainers. Even taking the favorable frame at face value,
the curated channel's advantage is concentrated in the explicit-danger
rules; on default-insecure practices (automount, limits, digest pinning)
charts are statistically indistinguishable from --- and on CPU limits
and digest pinning worse than --- production GitHub code. Curation
polishes the visible edge of the configuration surface; it does not
change the defaults underneath.

\hypertarget{rq2-precision-and-taxonomy-which-numbers-can-be-believed}{%
\subsubsection{4.3 RQ2 --- Precision and taxonomy: which numbers can be
believed}\label{rq2-precision-and-taxonomy-which-numbers-can-be-believed}}

Overall precision on the reconciled, author-adjudicated sample is
\textbf{0.633 {[}0.527, 0.714{]}} (279 TP / 162 FP over decided rows;
cluster-bootstrap interval by source), adjusting the 27,460 raw findings
to an estimated \textbf{$\approx$19,200 genuine misconfigurations} (95\% CI
{[}17,100, 23,100{]} --- a model-based extrapolation of the per-rule
estimates, not a count of individually verified findings). A
distinct-finding sensitivity analysis gives 0.626 --- the intentional
duplicate rows do not move the result. Population-weighted, overall
precision is 0.699 (Section 3.5); the per-rule adjustments that follow
are the same under either aggregate. But the aggregate conceals the
paper's most useful variation: precision is a property of the
\emph{rule}, and it splits the taxonomy into rules that can be read at
face value and rules that cannot.

\textbf{Face-value rules (precision $\approx$ 1.0).} Missing seccomp profile,
privilege escalation not disabled, root not prevented, and
\texttt{:latest} images are essentially never benign in context (all
1.000 on their samples); SA-token automount (0.857) and dangerous
capabilities (0.824, under the addendum's infrastructure test) are close
behind. For these, scanner output \emph{is} prevalence, up to sampling
error.

\textbf{Context-dominated rules (precision $\leq$ 0.4).} Committed Secret
objects are the extreme: \textbf{0.125}. Under the credential-material
test, most committed Secrets are empty, operator-populated, or
documented placeholders; the genuine exposure is a credential-bearing
minority that no content-blind rule can separate from the noise. Host
network (0.250) and host PID (0.214) findings are mostly node
infrastructure --- CNI/CSI agents and node exporters whose function
requires the access. Missing liveness probes (0.214) and the resource
rules (0.375--0.47) are dominated by production scoping: examples,
fixtures, and one-shot Jobs that a context-free count grades as
production gaps. Added capabilities (0.235) are mostly least-privilege
patterns (default-set additions, or \texttt{NET\_BIND\_SERVICE} over a
drop-ALL baseline).

\textbf{Table 2.} Per-rule precision (author-verified, reconciled to
engine v3; Wilson 95\% CIs). Selected rules; full table in
\texttt{results/precision\_provisional.md}.

\begin{longtable}[]{@{}
  >{\raggedright\arraybackslash}p{(\columnwidth - 6\tabcolsep) * \real{0.2308}}
  >{\raggedleft\arraybackslash}p{(\columnwidth - 6\tabcolsep) * \real{0.3077}}
  >{\raggedright\arraybackslash}p{(\columnwidth - 6\tabcolsep) * \real{0.2308}}
  >{\raggedright\arraybackslash}p{(\columnwidth - 6\tabcolsep) * \real{0.2308}}@{}}
\toprule\noalign{}
\begin{minipage}[b]{\linewidth}\raggedright
Rule
\end{minipage} & \begin{minipage}[b]{\linewidth}\raggedleft
Precision
\end{minipage} & \begin{minipage}[b]{\linewidth}\raggedright
95\% CI
\end{minipage} & \begin{minipage}[b]{\linewidth}\raggedright
Reading
\end{minipage} \\
\midrule\noalign{}
\endhead
\bottomrule\noalign{}
\endlastfoot
SEC\_\allowbreak{}NO\_\allowbreak{}SECCOMP & 1.000 & {[}0.82, 1.00{]} & face value \\
SEC\_\allowbreak{}APE\_\allowbreak{}NOT\_\allowbreak{}DISABLED & 1.000 & {[}0.80, 1.00{]} & face value \\
SEC\_\allowbreak{}ROOT\_\allowbreak{}NOT\_\allowbreak{}PREVENTED & 1.000 & {[}0.78, 1.00{]} & face value \\
IMG\_\allowbreak{}LATEST & 1.000 & {[}0.80, 1.00{]} & face value \\
SA\_\allowbreak{}TOKEN\_\allowbreak{}AUTOMOUNTED & 0.857 & {[}0.60, 0.96{]} & face value \\
SEC\_\allowbreak{}DANGEROUS\_\allowbreak{}CAPABILITY & 0.824 & {[}0.59, 0.94{]} & face value \\
RBAC\_\allowbreak{}WILDCARD\_\allowbreak{}VERBS & 0.733 & {[}0.48, 0.89{]} & mostly \\
RBAC\_\allowbreak{}WILDCARD\_\allowbreak{}RESOURCES & 0.647 & {[}0.41, 0.83{]} & mostly \\
SEC\_\allowbreak{}PRIVILEGED & 0.608 & {[}0.50, 0.71{]} & mixed \\
REL\_\allowbreak{}NO\_\allowbreak{}READINESS\_\allowbreak{}PROBE & 0.538 & {[}0.29, 0.77{]} & context-heavy \\
RES\_\allowbreak{}NO\_\allowbreak{}REQUESTS & 0.471 & {[}0.26, 0.69{]} & context-heavy \\
RES\_\allowbreak{}NO\_\allowbreak{}CPU\_\allowbreak{}LIMIT & 0.429 & {[}0.21, 0.67{]} & context-heavy \\
SECRET\_\allowbreak{}LITERAL\_\allowbreak{}ENV & 0.400 & {[}0.20, 0.64{]} & context-heavy \\
RES\_\allowbreak{}NO\_\allowbreak{}MEM\_\allowbreak{}LIMIT & 0.375 & {[}0.18, 0.61{]} & context-heavy \\
HOST\_\allowbreak{}NETWORK & 0.250 & {[}0.10, 0.49{]} & context-dominated \\
SEC\_\allowbreak{}CAPABILITY\_\allowbreak{}ADDED & 0.235 & {[}0.10, 0.47{]} & context-dominated \\
HOST\_\allowbreak{}PID / REL\_\allowbreak{}NO\_\allowbreak{}LIVENESS\_\allowbreak{}PROBE & 0.214 & {[}0.08, 0.48{]} &
context-dominated \\
SECRET\_\allowbreak{}COMMITTED & 0.125 & {[}0.03, 0.36{]} & context-dominated \\
\end{longtable}

\textbf{Taxonomy v2.} The validation exercise reorganizes the 25 rules
into a three-tier taxonomy, and the false positives into five reusable
context classes:

\begin{itemize}
\tightlist
\item
  \textbf{Tier 1 --- Misconfigurations.} States that are wrong in
  essentially all contexts (the face-value rules above, plus wildcard
  RBAC and cluster-admin binding at 0.65--0.80 with the own-group
  exception below).
\item
  \textbf{Tier 2 --- Hardening gaps.} States that are genuine but lesser
  gaps, reported separately and never folded into misconfiguration
  precision: digest-unpinned images \emph{with} an explicit version tag
  (7 of 14 sampled no-digest rows), and the context-heavy resource/probe
  rules when they occur in production code.
\item
  \textbf{Tier 3 --- Manifest rot.} Deprecated API versions: real,
  mechanically certain, and not a security misconfiguration; reported as
  their own line.
\end{itemize}

False positives concentrate in five context classes, each of which
became an explicit protocol test: \textbf{infra-need} (node agents whose
function requires the access, gated on a drop-ALL baseline for
capabilities), \textbf{placeholder} (Secrets and env values without
credential material), \textbf{fixture} (test/example content, decided by
content over source tag), \textbf{standard operator pattern} (wildcard
RBAC confined to the operator's own API groups; least-privilege
capability additions), and \textbf{not-exploitable} (root/APE findings
with no escalation path: non-root enforced and no added capabilities). A
scanner --- or a study --- that implements these five distinctions
recovers most of the precision the naive rules lose.

\hypertarget{rq3-auditing-the-auditors}{%
\subsubsection{4.4 RQ3 --- Auditing the
auditors}\label{rq3-auditing-the-auditors}}

On identical Helm-arm input with default configurations, approximate
coverage of the 25 practices is: Checkov \textasciitilde16, Polaris
\textasciitilde13, kube-linter \textasciitilde11, kubeaudit
\textasciitilde11 (partial coverage counted as one half). Three results
matter more than the ranking.

\textbf{Nobody covers the long tail, and one gap is total.} Digest
pinning --- the single most prevalent practice in the corpus --- is
checked by exactly one tool (Checkov CKV\_K8S\_43). Wildcard RBAC has no
\emph{default-enabled} check that catches it broadly: Checkov's
CKV\_K8S\_49 flags 29 of the 155 wildcard findings on this input, and
kube-linter ships wildcard checks that are off by default. And
\textbf{no tool, in the configuration run, inspects committed Secret
contents} (Checkov ships a separate secrets framework, outside the
Kubernetes checks we ran): all four flag (or ignore) Secret objects
without distinguishing credential material from placeholders, which
Section 4.3 shows is a 7:1 noise-to-signal distinction. The most
sensitive objects in the corpus are the ones the tooling understands
least.

\textbf{Agreement exists only where danger is explicit.} On privileged
containers, the four tools report 26--27 findings against our 27 ---
near-perfect consensus. On missing CPU limits, the same input yields
\textbf{343 / 411 / 416 / 497} findings across the tools against 502
container-level cases in our engine. The spread is bookkeeping, not
detection: kubeaudit counts at pod level and merges ``no limits at all''
(421) with ``limits without CPU'' (76); kube-linter merges limits with
requests (416); Checkov counts a narrower property (343). A consumer
comparing tools --- or tracking one tool over time --- sees a
\textasciitilde1.4$\times$ swing that has nothing to do with the underlying
manifests.

\textbf{Accuracy against verified ground truth.} Coverage and counts say
what a tool \emph{can} check and how much it \emph{reports}; the
verified sample lets us ask whether it is \emph{right}. Where a tool has
a check for a sampled practice, it detects verified true positives
reliably --- kubeaudit 0.98, Polaris 0.93, kube-linter 0.90 --- with
Checkov the outlier at 0.76 (its run-as-non-root check systematically
misses containers that run as root \emph{explicitly}: 8 of its 15
misses). Counting practices with no check as missed, effective detection
of genuine misconfigurations is 0.52--0.75, and the ranking inverts the
raw-count intuition: Polaris leads (0.75) on the strength of coverage,
not per-check detection. The symmetric result is the more damning one:
the same checks flag \textbf{86--95\% of the findings verification
judged false positives}. The taxonomy's context classes ---
infrastructure agents, standard operator patterns, least-privilege
capability sets --- fool the production tools at about the same rate
they fooled our unvalidated rules. The tools' deficit on real manifests
is not detection; it is coverage and context. These are sample-based
estimates from the 126 Helm-arm validation rows, drawn from the engine's
own findings: the flagging figures are shares of sampled findings that
verification judged benign, not false-positive rates over all negative
configurations, which a findings-only sample cannot measure. The
benchmark for ``benign'' is likewise ours: the tools implement the
definitional reading of each practice and do not claim the v3 context
classes, so the flagging rates measure the distance between definitional
checking and context-aware judgment --- the quantity a user who trusts
the output experiences --- rather than vendor error against the tools'
own specifications.

\begin{longtable}[]{@{}
  >{\raggedright\arraybackslash}p{(\columnwidth - 6\tabcolsep) * \real{0.2000}}
  >{\raggedleft\arraybackslash}p{(\columnwidth - 6\tabcolsep) * \real{0.2667}}
  >{\raggedleft\arraybackslash}p{(\columnwidth - 6\tabcolsep) * \real{0.2667}}
  >{\raggedleft\arraybackslash}p{(\columnwidth - 6\tabcolsep) * \real{0.2667}}@{}}
\toprule\noalign{}
\begin{minipage}[b]{\linewidth}\raggedright
Tool
\end{minipage} & \begin{minipage}[b]{\linewidth}\raggedleft
TP detection (check exists)
\end{minipage} & \begin{minipage}[b]{\linewidth}\raggedleft
TP detection (effective)
\end{minipage} & \begin{minipage}[b]{\linewidth}\raggedleft
FP-flagging (check exists)
\end{minipage} \\
\midrule\noalign{}
\endhead
\bottomrule\noalign{}
\endlastfoot
kube-linter & 0.90 {[}0.78, 0.96{]} & 0.52 & 0.92 \\
kubeaudit & 0.98 {[}0.87, 1.00{]} & 0.55 & 0.95 \\
Polaris & 0.93 {[}0.84, 0.97{]} & 0.75 & 0.87 \\
Checkov & 0.76 {[}0.64, 0.85{]} & 0.64 & 0.86 \\
\end{longtable}

\emph{Ground truth: 73 verified-TP and 50 verified-FP Helm-arm sample
rows; rates over rows where a mapped check exists, except the effective
column, which counts no-check practices as missed. Method, joins, and
caveats: \texttt{results/tool\_accuracy.md}.}

\textbf{Defaults decide what gets found.} kube-linter's probe checks
validate probe \emph{port correctness}, not probe existence; its
SA-token and wildcard checks are absent or opt-in; kubeaudit's SA check
covers only the default service account. The pattern is consistent: each
tool's default configuration is tuned to the explicit-danger tier of the
taxonomy, and the default-insecure tier --- where Sections 4.1 and 4.3
show most genuine misconfiguration lives --- is covered partially,
optionally, or not at all.

\begin{center}\rule{0.5\linewidth}{0.5pt}\end{center}

\hypertarget{discussion}{%
\subsection{5. Discussion}\label{discussion}}

\textbf{The risk is the default, and the default is a design outcome.}
The practices with the highest prevalence --- automounted tokens, absent
limits, unpinned images, absent seccomp --- are precisely those where
doing nothing produces the finding. Manifest authors are not choosing
insecurity; they are accepting a platform posture. A scope note on what
the findings are: most are \emph{exposures and hardening gaps}, not
demonstrated vulnerabilities --- an automounted token's impact depends
on the permissions bound to its service account and on cluster context,
which a manifest alone does not determine. The taxonomy's tiers (Section
4.3) keep that distinction in view; the prevalence claims are about
configuration posture, not exploitability. Interventions aimed at author
behavior (linters in editors, review checklists) fight the last war: the
durable fix is at the platform and distribution layer --- admission-time
defaults (Pod Security Standards, policy engines with the checks
\emph{enabled}), chart templates that ship hardened security contexts,
and image pipelines that pin digests at build time. Our curation result
(Section 4.2) is the cautionary evidence: even the ecosystem's most
maintained artifacts only partially overcome the defaults.

\textbf{Scanner output needs a validity label.} Two scanners can
disagree by 1.4$\times$ on the same manifests without either being ``wrong''
--- they are counting different units and different properties under the
same English name. We suggest the field adopt three reporting norms:
state the counting unit (pod vs container), state whether merged checks
are decomposed, and state per-check precision on a public validation
corpus, as this study does for its own engine. A count without those
three qualifiers is not comparable across tools or across time.

\textbf{The five context classes are a checklist for tool authors.} Most
of the precision our naive rules lost is recoverable with distinctions
that are mechanical once named: join the ServiceAccount object before
grading automount; require credential material before calling a Secret
committed; treat a drop-ALL baseline plus function-required capabilities
on a node agent as infrastructure, not exposure; scope wildcard RBAC by
API-group ownership; and never grade fixtures as production. Several of
these are one-join or one-test changes. The tools' roadmaps, not their
detection cores, are the bottleneck.

\textbf{For practitioners}, the actionable reading of Tables 1 and 2 is
short: the face-value rules are a reliable triage list (fix seccomp,
privilege escalation, root prevention, and \texttt{:latest} first ---
the findings are real), SA tokens and resource limits are the largest
genuine exposure by volume, and committed-Secret alerts should be routed
to content inspection rather than counted.

\hypertarget{threats-to-validity}{%
\subsection{6. Threats to validity}\label{threats-to-validity}}

\textbf{Construct validity} rests on the labeling protocol. We mitigate
by \emph{disclosing} rather than concealing its evolution: three
protocol versions, a blinded author--AI agreement subset ($\kappa$ = 0.526
against the v2-standard labels, with the disagreement decomposition
showing 4 of 6 tracing to the v2$\rightarrow$v3 standards change itself),
adjudicated rule-level amendments, and author adjudication of the labels
(Section 3.4 gives the row-level provenance). The ground truth is thus
single-author and AI-assisted; an independent human labeling under the
frozen protocol is in preparation and is the obvious strengthening step.
Readers who prefer the purely definitional standard can recompute
precision under it from the released sample (the v2 pre-labels are
preserved); the prevalence \emph{rankings} are unchanged under either
standard, though several per-rule precisions move, as Section 4.3
reports. The $\kappa$ figure is below the conventional 0.7 benchmark and we do
not claim otherwise; it measures agreement with labels produced under a
superseded standard, and the verification pass --- not the $\kappa$ subset ---
is the ground truth of record.

\textbf{Internal validity.} Engine v3's three semantic fixes (SA-object
join, init containers in probe rules, Unconfined seccomp) each moved
findings materially; residual engine bugs would bias prevalence and are
mitigated by the validation sample itself --- a systematically wrong
rule shows up as low precision, and none of the face-value rules do. For
RQ3, check-to-rule mapping was manual; we mitigated by reconciling
counts rule by rule and reporting merged checks as merged rather than
forcing attribution.

\textbf{External validity.} The GitHub arm over-represents public,
open-source practice and under-represents private enterprise manifests,
which plausibly differ in both directions (more policy enforcement, more
legacy). The Helm arm's best-chart bias is measured in Section 4.2 and
cuts toward \emph{understating} ecosystem-wide prevalence. All data is
from October 2026; defaults and tooling evolve, and the replication
package is designed for re-collection.

\textbf{Reliability.} Intra-rater consistency was 44/44 on hidden
duplicates. The sample-to-population reconciliation (451/457 rows
matched to engine v3) and the distinct-finding sensitivity analysis
(0.626 vs 0.633) bound the effect of the v2$\rightarrow$v3 frame change. Corpus
snapshots are pinned by commit SHA; the corpus is rebuildable from the
released collectors and logs.

\hypertarget{related-work}{%
\subsection{7. Related work}\label{related-work}}

\hypertarget{empirical-studies-of-kubernetes-misconfiguration}{%
\subsubsection{7.1 Empirical studies of Kubernetes
misconfiguration}\label{empirical-studies-of-kubernetes-misconfiguration}}

The closest predecessor is Rahman et al.'s empirical study of security
misconfigurations in open-source Kubernetes manifests
CITEKEY\_Rahman2023\_END. Mining 2,039 manifests from 92 repositories,
they characterize 11 categories of security misconfiguration, build a
detection tool (SLI-KUBE) to quantify their frequency (1,051
misconfigurations), and validate by reporting ten misconfigurations to
practitioners, who agreed to fix 60\%. We build on their
characterization in four ways: an order-of-magnitude larger corpus with
a second, curated arm (rendered Helm charts) that supports a curation
comparison; a broader practice set (25 rules across security, resources,
reliability, images, RBAC, and secrets, rather than security alone);
separation of production manifests from examples and fixtures, which
their manifest-level counting does not attempt; and a validation design
that estimates \emph{per-rule precision} over a stratified sample with a
documented protocol and a blinded agreement check, rather than
fix-agreement on a small set of reports. Their headline conclusion ---
insecure configurations are widespread --- is corroborated by our data
at larger scale; our contribution is establishing \emph{which} of the
counts can be believed, and by how much.

Minna et al.~mine Helm charts from Artifact Hub and compare seven
analyzer tools by their default policies, then evaluate whether LLMs can
mitigate the reported misconfigurations CITEKEY\_Minna2025\_END --- in
this journal. Their tool comparison aggregates common and unique
findings per tool and includes a manual false-positive check on a
subset. Our study asks the prior question their pipeline assumes away:
whether the tools' reports on natural input are \emph{correct}, rule by
rule, against verified ground truth --- and finds that for
context-sensitive rules they frequently are not (Section 4.3), a result
that also bears on their mitigation loop, since an LLM asked to fix a
false positive can only damage a working configuration. The studies are
complementary: they evaluate repair; we evaluate the measurement the
repair depends on. We also add the raw-manifest arm and the
production/fixture separation their chart-only corpus does not require.

A third recent line works from developer reports rather than manifests:
Ghorab et al.~derive a taxonomy of Kubernetes security misconfigurations
from 2,662 Stack Overflow issues, trace how the issues evolve with
project maturity, and evaluate LLM repair, reaching 98.5\% correction
accuracy with a schema-guided validation framework
CITEKEY\_Ghorab2026\_END. Their unit of analysis --- what developers
\emph{report struggling with} --- complements ours: what
production-shaped code \emph{contains}, and whether automated findings
about it are valid. Their repair results, like Minna et al.'s,
presuppose that reported misconfigurations are genuine, which is the
premise our validation tests.

Beyond Kubernetes, configuration errors have been studied empirically in
commercial and open-source systems CITEKEY\_Yin2011\_END, and the
systems approaches to tackling them surveyed in
CITEKEY\_XuZhou2015\_END; studies of infrastructure-as-code ecosystems
(Dockerfiles, Puppet, Terraform, CI configuration) generally find
insecure defaults and smell-like patterns to be widespread. Those
results motivate ours; the Kubernetes setting differs in that the
configuration surface is a single, structured API with platform-supplied
defaults for nearly every field, which is what makes a uniform 25-rule
measurement across two distribution channels possible.

\hypertarget{scanner-evaluation-and-benchmarks}{%
\subsubsection{7.2 Scanner evaluation and
benchmarks}\label{scanner-evaluation-and-benchmarks}}

Eom et al.~present the first systematic benchmark of Kubernetes
misconfiguration scanners CITEKEY\_Eom2026\_END: ten tools, whose
documented coverage they found too ambiguous to take at face value, so
they reverse-engineered precise detection targets from implementations
and dynamic testing, then measured recall on a constructed benchmark of
4,109 misconfiguration files (281 unique misconfigurations), finding
best recall of 43.8\% (static) and 62.3\% (dynamic). Their ground truth
is \emph{constructed} --- each benchmark file is a known
misconfiguration --- which is the right design for measuring recall, and
the wrong one for measuring what happens on real manifests, where the
open question is not ``does the tool fire on a seeded defect'' but ``is
this firing, on this real object, a genuine misconfiguration.'' That is
the quantity we measure: per-rule precision of findings on natural
corpora, against author-verified labels, plus prevalence --- neither of
which a seeded benchmark can report. The two results also converge from
opposite directions: their finding that scanners adopt different
criteria for the same documented misconfiguration is the
implementation-level explanation of our measurement-level finding that
identical input yields \textasciitilde1.4$\times$ count spreads through
counting semantics (Section 4.4). Where they audit ten scanners against
a specification, we audit four against the wild.

Industry comparisons of hardening guidelines and scanner scoring (e.g.,
the Dynatrace Research comparison of eight guidelines and ten scanners)
likewise report substantial coverage and scoring disparities; they share
our conclusion that scanner outputs are not mutually comparable, reached
there by guideline mapping rather than by validated findings.

\hypertarget{benchmarks-and-hardening-guidance}{%
\subsubsection{7.3 Benchmarks and hardening
guidance}\label{benchmarks-and-hardening-guidance}}

The CIS Kubernetes Benchmark and the NSA/CISA Kubernetes Hardening Guide
define the practice standard our rules operationalize, and the Pod
Security Standards define the platform's own posture tiers. Taxonomy v2
(Section 4.3) can be read as an empirical annotation of those documents:
which of their recommendations correspond to states that are wrong in
essentially every context (Tier 1), which to hardening opportunities
(Tier 2), and which to hygiene (Tier 3) --- and, through the five
false-positive context classes, where the guidance needs contextual
exceptions that the documents themselves do not state
(node-infrastructure agents being the largest).

\hypertarget{positioning}{%
\subsubsection{7.4 Positioning}\label{positioning}}

\begin{longtable}[]{@{}
  >{\raggedright\arraybackslash}p{(\columnwidth - 8\tabcolsep) * \real{0.2000}}
  >{\raggedright\arraybackslash}p{(\columnwidth - 8\tabcolsep) * \real{0.2000}}
  >{\raggedright\arraybackslash}p{(\columnwidth - 8\tabcolsep) * \real{0.2000}}
  >{\raggedright\arraybackslash}p{(\columnwidth - 8\tabcolsep) * \real{0.2000}}
  >{\raggedright\arraybackslash}p{(\columnwidth - 8\tabcolsep) * \real{0.2000}}@{}}
\toprule\noalign{}
\begin{minipage}[b]{\linewidth}\raggedright
Dimension
\end{minipage} & \begin{minipage}[b]{\linewidth}\raggedright
Rahman et al.~2023
\end{minipage} & \begin{minipage}[b]{\linewidth}\raggedright
Minna et al.~2025
\end{minipage} & \begin{minipage}[b]{\linewidth}\raggedright
Eom et al.~2026
\end{minipage} & \begin{minipage}[b]{\linewidth}\raggedright
This study
\end{minipage} \\
\midrule\noalign{}
\endhead
\bottomrule\noalign{}
\endlastfoot
Corpus & 2,039 manifests, 92 repos & Artifact Hub charts (all) & 4,109
constructed files & 16,243 docs, 2 arms (60 repos + 151 charts) \\
Practices & 11 security categories & tools' default policies & 281
misconfigurations (union of 10 tools) & 25 rules, 6 families \\
Ground truth & practitioner fix-agreement (10 reports) & manual FP check
on a subset & constructed by design & 457-finding stratified sample,
blinded subset, protocol + adjudication, author-adjudicated \\
Tool treatment & builds one tool & 7 tools, common/unique counts + LLM
repair & 10 tools, recall on benchmark & 4 tools on identical natural
input: coverage, count semantics, accuracy vs verified labels \\
Production vs fixture & not separated & n/a (charts) & n/a (constructed)
& separated; content-over-tag rule \\
Prevalence & manifest-level counts & chart-level counts & none
(benchmark) & workload-level, cluster-level, precision-adjusted \\
\end{longtable}

In one sentence: prior work establishes that Kubernetes misconfiguration
exists (Rahman et al.), that tools report it inconsistently and LLMs
repair it imperfectly (Minna et al.), and that scanners miss much of
what they claim to cover on seeded cases (Eom et al.). This study
establishes how much of it there is in production-shaped code, how much
of what any rule reports is real, and why the tools' numbers cannot be
compared naively --- the measurement foundation the other three
directions implicitly rely on.

\hypertarget{conclusion}{%
\subsection{8. Conclusion}\label{conclusion}}

Across 16,243 real-world documents, Kubernetes misconfiguration is not a
story of reckless operators but of permissive defaults accepted at
scale: tokens automounted in \textasciitilde9 of 10 production
workloads, limits and digest pinning absent as the norm, and the
explicitly dangerous settings comparatively rare. The measurement itself
needs care --- a third of naive findings are benign in context, and the
scanners the ecosystem relies on cover at most two-thirds of observed
practices, agree mainly where danger is explicit, and diverge by
bookkeeping where it is not. Both conclusions point the same way: harden
the defaults, and label the measurements. These claims describe public,
production-shaped manifests and rendered charts --- the configuration
the ecosystem publishes --- not the private state of running clusters,
which only their operators can observe.

\hypertarget{data-availability-and-replication-package}{%
\subsection{Data availability and replication
package}\label{data-availability-and-replication-package}}

All artifacts are public: corpus collectors and per-source logs, the
25-rule engine, per-arm and per-tag statistics, the RQ3 coverage matrix,
the complete labeling protocol (all versions and the 2026-10-27
addendum), the 457-row validation sample with evidence packets and final
labels, and the analysis scripts, at
https://github.com/mallyagirish/k8s-misconfig-study (Apache-2.0). GitHub
sources are pinned by commit SHA
(\texttt{results/github\_commit\_shas.json}) and Helm charts by name and
version (\texttt{results/helm\_render\_log.json}), so the corpus is
reconstructible as collected. An archival DOI (Zenodo) will be minted at
submission. The corpus manifests themselves remain under their original
licenses and are rebuilt, not redistributed.

\section*{Statements and Declarations}

\subsection*{Funding}
No funding was received for conducting this study.

\subsection*{Competing interests}
The author declares no competing interests.

\subsection*{Ethics approval and consent to participate}
Not applicable. The study analyzes publicly available open-source artifacts
and involves no human subjects.

\subsection*{Author contributions}
The author is the sole contributor to this work.

\subsection*{Use of artificial intelligence tools}
AI tools were used in two ways, both disclosed in the methods. First, an AI
assistant produced initial pre-labels for the validation sample under
labeling protocol v2; every label was subsequently verified or adjudicated
by the author under the governing protocol v3, and agreement was measured
on a blinded subset labeled by the author before seeing the pre-labels
(Section 3.4). Second, AI assistance was used for software development
support and language editing during preparation of this manuscript. The
author reviewed all output and takes full responsibility for the content.

\begin{thebibliography}{4}
\providecommand{\natexlab}[1]{#1}
\providecommand{\url}[1]{{#1}}
\providecommand{\urlprefix}{URL }
\providecommand{\doi}[1]{\url{https://doi.org/#1}}
\providecommand{\eprint}[2][]{\url{#2}}
 \bibcommenthead

\bibitem[{{Center for Internet Security}(2026)}]{CISBenchmark}
{Center for Internet Security} (2026) {CIS Kubernetes Benchmark}. Version
  current at time of study (accessed October 2026)

\bibitem[{{kube-linter, kubeaudit, Polaris, and Checkov
  projects}(2026)}]{ToolDocs}
{kube-linter, kubeaudit, Polaris, and Checkov projects} (2026) Tool
  documentation and check catalogues. Versions as recorded in the replication
  package

\bibitem[{{Kubernetes Project}(2026)}]{K8sPSS}
{Kubernetes Project} (2026) {Pod Security Standards} and security context
  documentation. Kubernetes.io/docs (accessed October 2026)

\bibitem[{{National Security Agency} and {Cybersecurity and Infrastructure
  Security Agency}(2022)}]{NSA2022}
{National Security Agency}, {Cybersecurity and Infrastructure Security Agency}
  (2022) {Kubernetes Hardening Guide}

\end{thebibliography}

\end{document}
